\begin{abstract}
確率的デコーダー（probabilistic decoder）の幾何は通常、そのヤコビ行列
（Jacobian）から読み取られる。
しかし、この方法ではブラックボックスデコーダーを扱えず、より根本的な逆問題
（inverse problem）、すなわち「出力分布間の局所距離から潜在曲面そのものを
決定できるか」という問いが残る。本論文は、明示的な凸多面体モデル
（convex polyhedral model）のもとで
この問いに肯定的な答えを与える。既知の球面型三角形分割上では、局所的な
2-Wasserstein距離が区分的ユークリッド（piecewise-Euclidean; PL）計量の辺長を定め、
同じ辺長から角欠損（angle defect）による厳密なPL曲率質量（curvature mass）が
得られる。それらが同一の三角形複体をもつ未知の狭義凸実現
（strictly convex realization）の辺であるとき、面の合同性（face congruence）と
剛性（rigidity）により、その実現は剛体運動（rigid motion）を除いて
$\mathbb R^3$ 内で同定される。デコーダーの
微分も3次元ラベルも用いない。

この復元過程のうち、計量復元と局所実現の部分を定量化する。形状正則性
（shape regularity）のもとで、加法的な辺長
誤差 $\varepsilon$ は頂点曲率質量を
$O(d_{\max}\varepsilon/\ell_{\min})$、スケール $h$ の準一様
（quasi-uniform）メッシュ上で面積正規化曲率を $O(\varepsilon/h^3)$ だけ摂動する。
相対辺長誤差はすべての内在距離（intrinsic distance）を制御し、したがって
Gromov--Hausdorff距離と
Gromov--Wasserstein距離も制御する。剛性行列（rigidity matrix）の条件により、
辺長誤差を剛体位置合わせ後の3次元頂点誤差へ局所的に変換でき、スペクトルギャップ
条件により球面上の距離から座標へ戻す段階も制御できる。
$s_{X_h}\gtrsim h^\beta$ を満たすメッシュ族では、得られるHausdorff誤差の
トレードオフは $O(\varepsilon_hh^{-(\beta+1)}+h^2)$ である。制御された
ベンチマークでは、全ランクかつ非可換なガウス共分散経路に単位球面を隠す。
デコーダー平均だけを
見ると球面は円板へ潰れる一方、局所Wasserstein距離はその計量全体を復元する。
標準球面に対する一連の実験は、安定な大域形状復元と、ノイズ感度の高い点ごとの曲率推定を
分離する。第2の辺長のみの実現実験は、最終段に球面制約を課さず、非球形楕円体からの
標本の凸包を復元する。
\end{abstract}

\section{はじめに}

平面地図は、それ自体では場所ごとの縮尺を保持しない。赤道付近と極付近で同じだけ
地図上を動いても、実際の移動距離は異なりうる。しかし、十分な数の局所移動時間が
分かれば、失われた縮尺、曲率、そして最終的には地球の形状を復元できる。本論文は、
生成デコーダーに対する同様の問題を扱う。潜在コードを地図座標、その座標でデコード
された確率分布を観測、局所的な最適輸送（optimal transport; OT）費用を測定された
移動時間とみなす。

既存研究の多くは、微分可能なデコーダーから出発し、そのヤコビ行列を用いて周囲計量
（ambient metric）を引き戻す（pullback）\cite{arvanitidis2018,chen2018,roy2022}。
この構成は
デコーダーが既知なら有用だが、有限個の距離クエリからどこまで幾何を決定できるかは
明らかにしない。また、計量テンソルや曲率ヒートマップは、まだ復元された曲面では
ないため、視覚的な出力にも曖昧さが残る。

そこで、本論文は次の明示的な逆問題を設定する。

\begin{quote}
三角形分割された潜在曲面（latent surface）と、その辺上でのみ観測された
ノイズ付き $W_2$ 距離から、内在曲率と隠れた3次元曲面をどの精度で復元できるか。
\end{quote}

入出力は意図的に有限とする。$S^2$ と同相な閉三角形複体を
$\mathcal T=(V,E,F)$ とする。共通の完備可分な基礎空間（ground space）を
$(\mathcal Y,d_{\mathcal Y})$ とし、$W_2$ の基礎費用を $d_{\mathcal Y}^2$ とする。
頂点 $i$ でデコーダーは $\mu_i\in\mathcal P_2(\mathcal Y)$ を返すが、推定器が受け取るのは
\begin{equation}
 \mathcal O=
 \left(\mathcal T,
 \{\widehat\ell_{ij}:(i,j)\in E\}\right),
 \qquad
 \widehat\ell_{ij}\simeq W_2(\mu_i,\mu_j).
 \label{eq:observations}
\end{equation}
のみである。$\mathbb R^3$ 内の正解点は評価時まで推定器に与えない。辺長から各面が
ユークリッド三角形として定まり、完備なPL計量が得られる。復元に現れる面角、曲率、面積、
内在距離はすべて、\eqref{eq:observations} だけの関数である。
以下の有限同定・安定性結果が用いるのは $E$ だけである。後の標準球面可視化では、
大域移動時間を近似するため3-hop局所問合せ集合
$E_{\mathrm q}\supset E$ を追加するが、その大きさは依然として線形であり、これらの
追加問合せは同一骨格定理には用いない。

外在的な復元では、観測後に周囲空間内の目標を定義するのではなく、別途次の仮定を
置く。

\begin{definition}[同一骨格凸観測モデル（same-skeleton convex observation model）]
\label{def:same-skeleton-model}
境界複体がちょうど $\mathcal T$ である未知の非退化な凸単体的多面体
（convex simplicial polyhedron）
$X_*=(x_i)_{i\in V}\subset\mathbb R^3$ が存在し、
\begin{equation}
 \ell_{ij}=W_2(\mu_i,\mu_j)=\norm{x_i-x_j}
 \qquad ((i,j)\in E).
 \label{eq:convex-observation-model}
\end{equation}
が成り立つ。すべての頂点が極点（extreme point）であり、すべての内二面角
（interior dihedral angle）が $\pi$ 未満である。このような実現を狭義凸実現
（strictly convex realization）と呼ぶ。測定値は上記の辺長を摂動したものである。
複体と凸モデルクラスは既知だが、$X_*$ は未知である。
\end{definition}

定義~\ref{def:same-skeleton-model} がなくても、辺観測は妥当なPL代理対象
$M_\ell$ を定めるが、既存の周囲曲面を復元するとは限らない。任意の端点間の
$W_2$ は周囲空間内の弦であって、内在測地距離ではない。
命題~\ref{prop:hidden-sphere} は、本論文のベンチマークについて
\eqref{eq:convex-observation-model} が厳密に成立することを示す。

\paragraph{本論文の貢献}
\begin{enumerate}
  \item 局所的なデコーダー分布間距離から、デコーダーのヤコビ行列を評価せずに
  PL潜在計量を定める\emph{Wasserstein移動時間復元}
  （Wasserstein travel-time reconstruction）を定式化する。
  \item 明示した同一骨格凸モデルのもとで、同定可能な三つの階層を接続する。
  辺長は角欠損曲率質量を定め、合同な三角形面は厳密な凸実現を定め、さらに
  非退化な剛性行列はその実現を局所的に安定にする。
  \item 辺ノイズから曲率、すべての内在距離、GH/GW差、および位置合わせ後の
  頂点誤差まで、明示的な摂動評価を与える。この評価は曲率のノイズ感度を示す。
  メッシュスケール $h$ では、加法的長さ誤差が面積正規化後に
  $\varepsilon/h^3$ へ増幅される。条件付き細分化の系により、剛性のスケーリングを
  明示的な解像度--ノイズ要件へ変換する。
  \item デコーダー平均には円板しか現れず、全ランクかつ非可換な共分散経路に欠落した高さを
  保存する正解既知のベンチマークを導入する。ガウス分布の $W_2$ 距離が内接球面多面体の辺長を
  厳密に復元するため、確率的観測モデルを厳密に検証できる。続く辺長のみの継続法実験
  により、楕円体状凸多面体（ellipsoidal convex polyhedron）でも同一骨格段階を検証する。
\end{enumerate}

\paragraph{OTに関する主張の範囲}
辺長が与えられた後の曲率、計量安定性、剛性に関する結果は、距離を入力とする主張であり、
最適輸送に固有ではない。一方、制御モデルはガウス $W_2$ のBures部分を実際に用いる。
共分散は全ランクで互いに非可換であり、生のFrobenius距離には
\eqref{eq:covariance-frobenius-distortion} の位置依存歪みがある。ただし、これは設計された
Wasserstein測地線であって、$W_2$ が唯一必要な距離であることの証明ではない。本論文が
主張する貢献は、問い合わせ可能なデコーダー分布間計量からPL曲率と凸逆復元までを結ぶ
一貫した定式化と保証であり、一般のグラフ実現法そのものの新規性ではない。

本論文は、位相を推定せず、一般の内在計量が一意な非凸埋め込みをもつとも主張せず、
MNISTに真の曲率が存在するとも主張しない。球面位相と凸性は仮定であり、滑らかな
球面は制御された評価対象である。

\section{関連研究}

\paragraph{潜在幾何（latent geometry）}
決定論的および確率的生成モデルの引き戻し計量は、デコーダーの微分を用いて潜在長と
測地線を定める\cite{arvanitidis2018,chen2018}。Optimal Latent Transportは、
デコーダー分布間のWasserstein計量を引き戻す\cite{roy2022}。Syrotaらは、座標が
同定不能であっても、深層潜在変数モデルの距離、角度、体積は同定可能となりうることを
示した\cite{syrota2025}。これらの研究は、既知の微分可能モデルに対して幾何を構成、
または利用する。本論文の対象は異なり、疎かつ微分を用いないWasserstein距離観測からの
復元精度を扱う。

\paragraph{距離から多様体へ（distances to manifolds）}
QiuとDaiは、射影された球面、タクシー移動時間、MNISTを含む例に対して、ノイズ付き
内在距離応答から滑らかなリーマン計量（Riemannian metric）とその微分を推定する
\cite{qiu2023}。
Feffermanらは、ノイズ付き内在距離から有界幾何（bounded geometry）をもつ抽象多様体を
有限標本保証付きで
復元する\cite{fefferman2020}。したがって、距離から計量を復元すること自体は本論文の
新規性ではない。本論文は観測を局所的なデコーダー $W_2$ クエリに限定し、内在計量から
離散曲率、さらに凸な外在実現（extrinsic realization）まで進む。Isomapも局所距離を
連結するが、その目的は
曲率や凸曲面の同定ではなく、ユークリッド座標表現を得ることである
\cite{tenenbaum2000}。
Hammらは、Wasserstein空間内の標本化された部分多様体について、漸近的なGWの意味で
内在計量と接空間を復元する\cite{hamm2025}。これに対し本論文は有限PL対象を採用し、
その曲率質量と同定可能な凸実現までを扱う。

\paragraph{ユークリッド距離幾何（Euclidean distance geometry）}
不完全な距離集合からユークリッド座標を復元する問題は古典的な距離幾何であり、
グラフ剛性、行列補完、センサ位置推定を含む\cite{liberti2014}。本論文の辺長のみの
実現もこの問題の一例であり、新しい一般グラフ実現法ではない。相違は、その前段の観測と
保証の連鎖にある。すなわち、局所的なデコーダー分布間距離からPL計量と曲率を定めた後、
明示した凸モデルのもとで外在実現へ進む。

\paragraph{多面体幾何と形状比較（polyhedral geometry and shape comparison）}
Alexandrovの定理は、内在PL計量によって凸多面体を特徴づける。BobenkoとIzmestievは
その構成的証明を与えた\cite{alexandrov2005,bobenko2008}。角欠損はPL曲面の内在曲率
測度である。滑らかなガウス曲率（Gaussian curvature）を近似するにはメッシュに
仮定が必要であり、任意の
三角形分割上で点ごとに整合的とは限らない\cite{gawlik2020,hartig2021}。そこで本論文は、
まず厳密なPL対象に対する有限ノイズ定理を述べ、滑らかな球面に対する誤差は別に報告する。
Gromov--Wasserstein距離は、測度保存等長写像（measure-preserving isometry）を除いて
距離測度空間を比較する
\cite{memoli2011}。既知の頂点対応を用いれば、その距離歪みから、より単純な上界も得られる。

\section{Wasserstein辺幾何}
\label{sec:problem}

各面 $f=(i,j,k)\in F$ について、正の長さ
$\ell_{ij},\ell_{jk},\ell_{ki}$ が狭義三角不等式を満たすとする。得られる
ユークリッド三角形を
等しい辺に沿って貼り合わせると、PL計量空間
$M_\ell=(\mathcal T,d_\ell)$ が定まる。この対象は内在的であり、その定義には
周囲ユークリッド空間内の座標を一切用いない。

$i$ を含む面における頂点 $i$ の角は、余弦定理から
\begin{equation}
 \theta_i^{jk}(\ell)
 =\arccos\!\left(
 \frac{\ell_{ij}^2+\ell_{ik}^2-\ell_{jk}^2}
 {2\ell_{ij}\ell_{ik}}
 \right).
 \label{eq:cosine-angle}
\end{equation}
と計算できる。内部頂点における厳密なPL曲率質量と、その重心面積密度
（barycentric-area density）は
\begin{align}
 \kappa_i(\ell)
 &=2\pi-\sum_{f\ni i}\theta_{i,f}(\ell),
 \label{eq:angle-defect}\\
 A_i(\ell)
 &=\frac13\sum_{f\ni i}A_f(\ell),
 &K_i(\ell)&=\frac{\kappa_i(\ell)}{A_i(\ell)}.
 \label{eq:discrete-density}
\end{align}
である。$\kappa_i$ はラジアン（radian）単位の積分曲率質量であって、点ごとの
ガウス曲率ではない。
正規化量 $K_i$ は滑らかな曲面との比較に有用だが、選択した頂点面積にも依存し、
任意のメッシュ上で点ごとの整合性をもつとは限らない。任意の閉球面三角形分割では、
ユークリッド三角形の内角和とオイラー（Euler）の公式から
\begin{equation}
 \sum_{i\in V}\kappa_i=4\pi
 \label{eq:discrete-gauss-bonnet}
\end{equation}
が成り立つ。

\subsection{ガウス型不確実性に隠れた球面}

次の構成が、実験で用いる制御された例である。また、平均だけに基づく距離では
不十分である理由も明確にする。

\begin{proposition}[非可換共分散に高さを隠すWasserstein恒等式]
\label{prop:hidden-sphere}
$R>0$ と $k\in(0,1)$ を固定し、
\begin{equation}
 B=\begin{pmatrix}1&0\\0&2\end{pmatrix},\qquad
 H=k\begin{pmatrix}0&1\\1&0\end{pmatrix},\qquad
 \lambda=\frac{4R^2}{3k^2}
 \label{eq:hidden-decoder}
\end{equation}
とする。$|x_3|\le R$ を満たす $x=(x_1,x_2,x_3)$ に対し、
\begin{equation}
 \begin{aligned}
  t_x&=\frac{x_3+R}{2R}, &m_x&=(x_1,x_2),\\
  C_x&=\lambda(I_2+t_xH)B(I_2+t_xH)
 \end{aligned}
 \label{eq:hidden-covariance-path}
\end{equation}
と置き、$\mu_x=\mathcal N(m_x,C_x)$ とする。すべての $C_x$ は正定値であり、
このスラブ内の任意の $x,y$ に対して、ガウスBures共分散距離 $d_{\mathrm B}$ は
\begin{equation}
 \begin{aligned}
 W_2^2(\mu_x,\mu_y)
 &=\norm{m_x-m_y}^2+d_{\mathrm B}^2(C_x,C_y)\\
 &=\norm{x-y}^2.
 \end{aligned}
 \label{eq:hidden-identity}
\end{equation}
が成り立つ。$x_3\ne y_3$ なら $C_xC_y\ne C_yC_x$ である。$S_R^2$ 上では、平均写像は
円板となり、同じ $(x_1,x_2)$ をもつ北半球と南半球の点を同一視するが、
Wasserstein観測はそれらを区別する。
\end{proposition}

\begin{proof}
$A_t=I_2+tH$ および $L_t=\sqrt\lambda A_tB^{1/2}$ と置けば
$C_t=L_tL_t^\trans$ である。$A_s$ と $A_t$ は $H$ の可換な正定値多項式なので、
$L_s^\trans L_t$ は正定値である。したがってガウスBures距離に対する直交Procrustes公式では
恒等変換が選ばれ、\cite{givensshortt1984,takatsu2011}
\begin{equation}
 \begin{aligned}
 d_{\mathrm B}^2(C_s,C_t)
 &=\norm{L_s-L_t}_F^2\\
 &=3\lambda k^2(s-t)^2
 =4R^2(s-t)^2
 \end{aligned}
 \label{eq:bures-height-identity}
\end{equation}
を得る。$s=t_x$、$t=t_y$ を代入すると
$d_{\mathrm B}(C_x,C_y)=|x_3-y_3|$ となり、\eqref{eq:hidden-identity} が従う。
$t_x\in[0,1]$ と $k<1$ から正定値性も分かる。直接計算により
\[
 C_sC_t-C_tC_s
 =3\lambda^2k(s-t)(1+k^2st)
 \begin{pmatrix}0&1\\-1&0\end{pmatrix}
\]
であり、$s\ne t$ なら非零である。一方、生の共分散座標は
\begin{equation}
 \norm{C_s-C_t}_F
 =\lambda k|s-t|\sqrt{18+5k^2(s+t)^2}
 \label{eq:covariance-frobenius-distortion}
\end{equation}
を満たすので、そのFrobenius距離を一つの定数で再尺度化しても高さ計量は復元できない。
\end{proof}

球面上の2点間の大円距離を $d\in[0,\pi R]$ とすると、
\eqref{eq:hidden-identity} のWasserstein距離は弦長
\begin{equation}
 W_2(\mu_x,\mu_y)=2R\sin\!\left(\frac{d}{2R}\right),
 \qquad
 0\le d-W_2\le \frac{d^3}{24R^2}.
 \label{eq:chord-arc}
\end{equation}
である。したがって、遠方の2点に対する単一の $W_2$ クエリは球面内部を近道する。
局所クエリを連結すれば、メッシュ細分化に伴ってこの近道は解消される。このため、
観測グラフは全点対ではなく局所辺から構成する。

\section{何が復元可能か}
\label{sec:theory}

厳密な同定可能性（identifiability）とノイズ下での安定性を分けて論じる。最初の結果は
古典的だが、「曲面を復元する」という表現に欠けていた条件を与える。

\begin{theorem}[厳密な同一骨格同定]
\label{thm:convex-identification}
定義~\ref{def:same-skeleton-model} のもとで、厳密な局所Wasserstein辺長は、
ユークリッド群（Euclidean group）の作用を除いて $X_*$ を決定する。
\end{theorem}

\begin{proof}
三辺相等（side--side--side; SSS）条件により、対応する各三角形面は合同である。
したがってコーシー（Cauchy）の剛性定理により、対応面が合同な二つの凸多面体は合同となる
\cite{alexandrov2005,pak2006}。
\end{proof}

この定理は、同一複体をもつ凸な正解対象の存在を条件とし、凸性や組合せ構造そのものを
推定するわけではない。Alexandrovの定理は、適切な正曲率PL計量に対して、これとは異なる
純粋に内在的な実現結果を与える\cite{bobenko2008}。しかし、その実現における自然な
多面体辺は、指定した入力グラフの辺と一致するとは限らない。したがって本論文では、
Alexandrovの定理と後述する辺剛性評価を合成しない。

\subsection{辺ノイズから曲率ノイズへ}

$d_i$ を頂点 $i$ に接続する面の数とし、$d_{\max}=\max_i d_i$ と置く。次の形状正則
クラスを用いる。真の三角形と
摂動後の三角形はいずれも、辺長が
$[\ell_{\min}/2,2\ell_{\max}]$ に含まれ、すべての角が $\alpha_0>0$ 以上であるとする。
また $\rho=\ell_{\max}/\ell_{\min}$ と置く。これらの条件により、任意に退化へ近づく
三角形を除外する。

\begin{theorem}[離散曲率の安定性]
\label{thm:curvature-noise}
次を仮定する。
\begin{equation}
 \max_{e\in E}\abs{\widehat\ell_e-\ell_e}
 \le\varepsilon.
 \label{eq:additive-noise}
\end{equation}
上記の形状正則性のもとで、定数
$C_\theta=C_\theta(\rho,\alpha_0)$ および
$C_A=C_A(\rho,\alpha_0)$ が存在し、
\begin{align}
 \abs{\widehat\kappa_i-\kappa_i}
 &\le C_\theta d_i
       \frac{\varepsilon}{\ell_{\min}},
 \label{eq:defect-stability}\\
 \abs{\widehat A_i-A_i}
 &\le C_A d_i\ell_{\max}\varepsilon.
 \label{eq:area-stability}
\end{align}
を満たす。\eqref{eq:area-stability} の右辺が $A_i/2$ 以下なら、
\begin{equation}
 \abs{\widehat K_i-K_i}
 \le
 \frac{2C_\theta d_i\varepsilon}
      {A_i\ell_{\min}}
 +\frac{2\abs{\kappa_i}C_A d_i\ell_{\max}\varepsilon}
      {A_i^2}.
 \label{eq:density-stability}
\end{equation}
特に、
\begin{equation}
 C_\theta d_{\max}\frac{\varepsilon}{\ell_{\min}}
 <\kappa_{\min}:=\min_i\kappa_i.
 \label{eq:convexity-noise}
\end{equation}
なら、正の角欠損質量は摂動後も保たれる。
\end{theorem}

\begin{proof}
辺長三つ組からなるコンパクトな形状正則集合上で余弦定理
\eqref{eq:cosine-angle} を微分すると、各角度勾配の $\ell_1$ ノルムは
$C_\theta/\ell_{\min}$ で抑えられる。$1/\sin\theta$ の因子は $\alpha_0$ により
制御される。平均値の定理を用い、頂点に接続する面について和を取れば
\eqref{eq:defect-stability} を得る。同じコンパクト集合上でヘロン（Heron）の公式は滑らかで、
その勾配は $C_A\ell_{\max}$ で抑えられるため、\eqref{eq:area-stability} が従う。
最後に、
\[
 \frac{\widehat\kappa_i}{\widehat A_i}
 -\frac{\kappa_i}{A_i}
 =\frac{\widehat\kappa_i-\kappa_i}{\widehat A_i}
 +\kappa_i\frac{A_i-\widehat A_i}{A_i\widehat A_i},
\]
および $\widehat A_i\ge A_i/2$ から \eqref{eq:density-stability} を得る。
条件~\eqref{eq:convexity-noise} により、摂動後の角欠損はすべて正となる。
\end{proof}

辺スケールが $h$、次数が有界、かつ $A_i\asymp h^2$ である準一様メッシュでは、
定理~\ref{thm:curvature-noise} は
\begin{equation}
 \abs{\widehat\kappa_i-\kappa_i}
 =O(\varepsilon/h),
 \qquad
 \abs{\widehat K_i-K_i}=O(\varepsilon/h^3).
 \label{eq:noise-amplification}
\end{equation}
となる。前者は厳密なPL曲率測度である。後者は縮小する面積で除すため、必然的に
ノイズ感度が高い。式~\eqref{eq:noise-amplification} は摂動評価であって、任意の
メッシュ上で滑らかな曲率へ点ごとに収束するとの主張ではない。辺誤差が相対的で
$\varepsilon\asymp\eta h$ なら、二つのオーダーはそれぞれ $O(\eta)$ と
$O(\eta/h^2)$ になる。さらに、メッシュ族が $\kappa_i\asymp h^2$ という意味で
曲率整合的（curvature-consistent）であるとき、十分小さい定数を伴う
$\eta=O(h^2)$ は、すべての角欠損の符号を保つための十分条件となる。
したがって、固定した相対ノイズは、細分化のもとで大域距離を保存しながら点ごとの
曲率情報を失わせうる。実験では両方の効果を報告する。

\subsection{内在誤差と外在誤差}

相対長さノイズからは、スケールに依存しない保証が得られる。対応する真の面と摂動後の
面の各組には、重心座標（barycentric coordinates）を保つ標準的なアフィン写像がある。
形状正則性により、その特異値は安定となる。

\begin{theorem}[距離および距離空間の安定性]
\label{thm:distance-stability}
形状正則性と、十分小さい $\eta$ に対する
\begin{equation}
 \max_{e\in E}
 \abs{\widehat\ell_e/\ell_e-1}\le\eta
 \label{eq:relative-noise}
\end{equation}
を仮定する。このとき $C_g=C_g(\rho,\alpha_0)$ が存在し、面ごとのアフィン対応
$\Phi:M_\ell\to M_{\widehat\ell}$ は
\begin{equation}
 (1-C_g\eta)d_\ell(p,q)
 \le d_{\widehat\ell}(\Phi p,\Phi q)
 \le(1+C_g\eta)d_\ell(p,q).
 \label{eq:all-distance-stability}
\end{equation}
を満たす。したがって、
\begin{equation}
 d_{\mathrm{GH}}(M_\ell,M_{\widehat\ell})
 \le\frac{C_g\eta}{2}\diam(M_\ell).
 \label{eq:gh-bound}
\end{equation}
さらに、固定された一様頂点重みを含む $M_\ell$ 上の任意の確率測度 $\nu$ に対し、
$p\mapsto(p,\Phi p)$ が誘導するカップリング（coupling）から
\begin{equation}
 GW_2\!\left((M_\ell,d_\ell,\nu),
 (M_{\widehat\ell},d_{\widehat\ell},\Phi_\#\nu)\right)
 \le C_g\eta\diam(M_\ell).
 \label{eq:gw-bound}
\end{equation}
を得る。
\end{theorem}

\begin{proof}[証明の概略]
ユークリッド三角形のグラム行列（Gram matrix）は、三つの辺長の二乗の滑らかな関数である。
したがって、
コンパクトな形状正則クラス上では、\eqref{eq:relative-noise} によりアフィン写像の
特異値を $1\pm C_g\eta$ で抑えられる。各可長曲線を面ごとに写し、長さ評価を積分して
下限を取ることで \eqref{eq:all-distance-stability} を得る。$\Phi$ のグラフは、歪みが
高々 $C_g\eta\diam(M_\ell)$ の対応（correspondence）であるため、
\eqref{eq:gh-bound} が従う。$GW_2$ の定義で誘導された測度カップリングを用いれば、
最後の主張を得る\cite{memoli2011}。
\end{proof}

式~\eqref{eq:gw-bound} では、意図的に二つの計量上で同じ抽象確率測度を用いている。
各計量について別々に正規化した面積測度は、互いの厳密な押し出し（pushforward）では
ない。これらの自然な測度を比較するには、密度摂動に由来する追加項が必要である。

球面実験にはもう一つの写像がある。近似された全球距離をスペクトル分解によって
3次元座標へ変換する写像である。次の命題は、この最終段の安定性だけを切り出す。

\begin{proposition}[球面スペクトル復元の局所安定性]
\label{prop:spherical-spectral-stability}
$X\in\R^{n\times3}$ の各行が $x_i\in S_R^2$ を満たすとし、$G=XX^\trans$、
$\gamma:=\lambda_3(G)>0$ とする。このとき $\lambda_4(G)=0$ である。さらに
\begin{equation}
 D_{ij}=R\arccos\!\left(\frac{x_i^\trans x_j}{R^2}\right)
 \label{eq:true-spherical-distance-matrix}
\end{equation}
と定める。対称行列 $\widehat D$ と $\widehat R>0$ に対し、
$\widehat G_{ij}=\widehat R^2\cos(\widehat D_{ij}/\widehat R)$ を作る。
その最大の三つの固有値を残し、得られる半正定値行列を
$\widehat Y\widehat Y^\trans$ と因数分解し、行ごとに
$\widehat x_i=\widehat R\widehat y_i/\norm{\widehat y_i}$ とする。
$(D,R)$ の近傍と有限な $C_{\mathrm{sph}}=C_{\mathrm{sph}}(X)$ が存在し、その近傍では
残した固有値は正で全行は非零となり、
\begin{equation}
 \min_{Q\in O(3)}\norm{\widehat X-XQ}_F
 \le C_{\mathrm{sph}}\left(
 \norm{\widehat D-D}_F+\sqrt n\abs{\widehat R-R}
 \right)
 \label{eq:spherical-spectral-stability}
\end{equation}
が成り立つ。上位三つの固有値相互の分離は仮定しない。
\end{proposition}

\begin{proof}[証明の概略]
$B(D,r)_{ij}=r^2\cos(D_{ij}/r)$ と置く。$d$ と $r$ に関するスカラー微分はそれぞれ
$-r\sin(d/r)$ と $2r\cos(d/r)+d\sin(d/r)$ であり、$(D,R)$ の近傍で有界である。
したがって $B$ は \eqref{eq:spherical-spectral-stability} のノルムに関して局所Lipschitz
である。Weylの不等式と $\gamma>0$ により、上位3次元固有空間は残りのスペクトルから
分離される。この孤立固有値クラスタへの射影はRieszスペクトル射影であり、ギャップが
開いている限り滑らかに変化するので、ランク3の正のスペクトル切断も局所Lipschitzとなる。
完全性のため、二つの列フルランク因子を直交Procrustesで位置合わせする。得られる水平
スライス上で、$Y\mapsto YY^\trans$ の $X$ における微分は
$E\mapsto XE^\trans+EX^\trans$ である。$X$ の特異ベクトル基底における直接のブロック計算
から、その利得は下から $\sqrt{2}\,\sigma_3(X)=\sqrt{2\gamma}$ で抑えられ、二次剰余
$EE^\trans$ は局所的に吸収できる。したがってグラム因子の摂動評価として、局所的に
\[
 \min_{Q\in O(3)}\norm{\widehat Y-XQ}_F
 \le \frac{C_0}{\sqrt\gamma}\norm{\widehat G-G}_F
\]
を得る。この評価は上位固有空間内部で固有値が重複しても成立する。最後に、
$\norm x=R$ と $y\ne0$ に対して
$\|\widehat R y/\|y\|-x\|\le|\widehat R-R|+2\|y-x\|$ である。
これを行ごとに適用すれば結論を得る。
\end{proof}

この命題は $\widehat D$ と $\widehat R$ を構成した後から始まる。制御ベンチマークで用いる
グラフ最短路、直径補正、半径推定そのものを保証するものではない。

定理~\ref{thm:distance-stability} は内在的な結果である。復元した3次元頂点を測るため、
$X\in\mathbb R^{|V|\times3}$ を狭義凸な三角形分割実現とし、その辺長写像を
$\mathcal L(X)=(\norm{X_i-X_j})_{(i,j)\in E}$ とする。このヤコビ行列が剛性行列
$R_X$ である。6次元の無限小ユークリッド運動の空間を
$\mathcal Z_X=\ker R_X$ と書き、座標上のFrobeniusノルムに関して
\begin{equation}
 s_X=\sigma_{\min}\!\left(R_X\big|_{\mathcal Z_X^\perp}\right)
 \label{eq:rigidity-singular-value}
\end{equation}
と定める。

\begin{proposition}[局所的な同一骨格3次元安定性]
\label{prop:rigidity}
狭義凸三角形分割実現では $s_X>0$ である。$\mathcal L(X)$ に十分近い任意の辺長ベクトル
$\widehat\ell$ は、ユークリッド運動を除いて一意な、複体 $\mathcal T$ をもつ近傍の狭義凸実現
$\widehat X$ をもち、
\begin{equation}
 \min_{Q\in O(3),\,b\in\R^3}
 \norm{\widehat X-(XQ+\mathbf 1b^{\trans})}_{F}
 \le
 \frac{2}{s_X}
 \norm{\widehat\ell-\mathcal L(X)}_2.
 \label{eq:rigidity-bound}
\end{equation}
を満たす。辺誤差が0へ近づくとき、係数 $2$ は $1+o(1)$ に置き換えられる。
\end{proposition}

\begin{proof}[証明の概略]
凸三角形分割多面体は無限小剛性（infinitesimal rigidity）をもつ\cite{pak2006}。
球面三角形分割では $|E|=3|V|-6$ である。したがって、$X$ を通り
$\mathcal Z_X^\perp$ に接するアフィンスライスを取ると、$\mathcal L$ の微分行列は正方かつ
非特異となり、その最小特異値は $s_X$ である。逆関数定理により局所実現と一意性を得る。
近傍を十分縮小すれば逆写像の微分ノルムを $2/s_X$ 以下にでき、狭義凸性も保たれる。
剛体運動について最小化すればスライス上の変位は増えないので、
\eqref{eq:rigidity-bound} が従う。この結果は局所的であり、制約のない最適化法がこの
分枝を見つけるとは主張しない。
\end{proof}

定理~\ref{thm:convex-identification} と命題~\ref{prop:rigidity} を合わせると、厳密な
大域同定と局所的な数値条件を区別できる。定数 $s_X$ は各実現に固有であり、細分化の
もとで0から一様に離れているとは仮定しない。

\subsection{滑らかな球面は別の近似対象である}

辺長を厳密に復元して同定されるのは有限多面体であり、有限解像度での滑らかな球面では
ない。残る誤差は、観測ノイズとは独立に評価できる。

\begin{proposition}[内接球面多面体による近似]
\label{prop:inscribed-sphere}
頂点 $x_i$ が半径 $R$ の球面上にあり、その凸包 $P_h$ が原点を含み、
$\mathcal T$ がその単体的境界であるとする。もし
\begin{equation}
 h=\max_{(i,j)\in E}d_{S_R^2}(x_i,x_j)<\frac{\pi R}{2},
 \label{eq:spherical-mesh-size}
\end{equation}
なら、命題~\ref{prop:hidden-sphere} の厳密な観測は $\partial P_h$ の内在計量を復元する。
さらに、
\begin{equation}
 d_{\mathrm H}(\partial P_h,S_R^2)
 \le R\bigl(1-\cos(h/R)\bigr)
 \le\frac{h^2}{2R}.
 \label{eq:inscribed-sphere-bound}
\end{equation}
が成り立つ。ノイズ付き辺長が命題~\ref{prop:rigidity} の局所近傍にとどまるとき、
$\widehat P$ を頂点 $\widehat X$ をもつ凸多面体とする。このとき、剛体位置合わせ後に
\begin{equation}
 d_{\mathrm H}(\partial\widehat P,S_R^2)
 \le
 \frac{2}{s_X}\norm{\widehat\ell-\ell}_2
 +R\bigl(1-\cos(h/R)\bigr)
 \label{eq:noisy-sphere-decomposition}
\end{equation}
を得る。本命題のHausdorff距離はすべて周囲ユークリッド距離を用いる。
\end{proposition}

\begin{proof}
式~\eqref{eq:hidden-identity} により、復元される各面は $P_h$ の対応面と合同になる。
定理~\ref{thm:convex-identification} は同一複体をもつ凸実現を同定する。$x_i$ を含む面について
$u=x_i/R$ と置く。この面の各頂点 $x_j$ は
$u^\trans x_j\ge R\cos(h/R)$ を満たす。したがって、面内の任意の凸結合 $q$ について
\[
 \norm q\ge u^\trans q\ge R\cos(h/R).
\]
また、$q$ は凸包に属するので $\norm q\le R$ である。したがって放射射影
（radial projection）が $q$ を動かす距離は高々 $R(1-\cos(h/R))$ である。
逆に、原点が $P_h$ に含まれるため、原点から出る任意の半直線は
$\partial P_h$ と交わる。その交点に同じ評価を適用すると、逆向きの有向Hausdorff評価も得られる。
初等不等式 $1-\cos t\le t^2/2$ から \eqref{eq:inscribed-sphere-bound} を得る。
命題~\ref{prop:rigidity} と三角不等式により \eqref{eq:noisy-sphere-decomposition} が従う。
\end{proof}

特に、一様誤差の仮定
$\max_e|\widehat\ell_e-\ell_e|\le\varepsilon$ から
\begin{equation}
 d_{\mathrm H}(\partial\widehat P,S_R^2)
 \le
 \frac{2\sqrt{|E|}}{s_X}\,\varepsilon
 +\frac{h^2}{2R}.
 \label{eq:end-to-end-sphere}
\end{equation}
を得る。式~\eqref{eq:end-to-end-sphere} は、「復元した多様体がどれだけ近いか」という
問いに対する本論文の最も単純な答えである。一方の項が局所Wasserstein観測誤差を、
もう一方が有限メッシュ近似誤差を測る。

\begin{corollary}[条件付き解像度--ノイズのトレードオフ]
\label{cor:conditional-refinement}
固定された球面上の準一様メッシュ族について、$|E_h|\le C_E h^{-2}$ および
$s_{X_h}\ge c_s h^\beta$（$\beta\ge0$）を仮定する。各解像度で復元が
命題~\ref{prop:rigidity} の局所分枝内にとどまるとし、その凸多面体を
$\widehat P_h$ と書く。このとき
\begin{equation}
 d_{\mathrm H}(\partial\widehat P_h,S_R^2)
 \le \frac{2\sqrt{C_E}}{c_s}\,
 \varepsilon_hh^{-(\beta+1)}+\frac{h^2}{2R}
 \label{eq:conditional-refinement-rate}
\end{equation}
が成り立つ。したがって $\varepsilon_h=o(h^{\beta+1})$ は整合性の十分条件であり、
$\varepsilon_h=O(h^{\beta+3})$ なら $O(h^2)$ の離散化率を保つ。Hausdorff誤差
$\delta$ を目標とするには、$h\asymp\delta^{1/2}$、
$\varepsilon_h=O(\delta^{(\beta+3)/2})$ とすれば十分であり、必要な辺数は
$O(\delta^{-1})$ である。同様に、与えられた小さな $\varepsilon$ に対して二項を
釣り合わせると、誤差は $O(\varepsilon^{2/(\beta+3)})$ となる。
\end{corollary}

\begin{proof}
$\sqrt{|E_h|}\le\sqrt{C_E}h^{-1}$ と
$s_{X_h}^{-1}\le c_s^{-1}h^{-\beta}$ を \eqref{eq:end-to-end-sphere} に代入する。
残りの主張は二項を比較、または釣り合わせれば従う。
\end{proof}

式~\eqref{eq:noisy-sphere-decomposition} の二項は意味が異なる。第一項は有限PL対象に
対する観測・復元誤差、第二項はその対象と滑らかな球面の間の離散化誤差である。
本論文では両者を分けて報告し、ノイズが0の多面体を厳密な滑らかな球面として提示しない。
剛性定数は固定された多面体に局所的な量であり、細分化のもとで一様に保たれるとは
主張しない。
